\section{Erd\H{o}s Problem \#481: collisions under expanding affine maps}
\noindent Complete harmonic-weight reconstruction, 24 September 2026.

\subsection*{Statement}
Let $a_i,b_i\in\mathbb N$ for $1\le i\le r$, with $C:=\sum_i1/a_i>1$. Starting from the one-element sequence $A_1=(1)$, replace each entry $x$ at every step by all $a_ix+b_i$, retaining multiplicities. We prove that some level $A_k$ has repeated entries. In particular, collisions at different levels alone would not suffice.

\subsection*{Harmonic mass cannot grow both linearly and exponentially}
Assume every level has distinct entries, and put
\[S_k=\sum_{x\in A_k}\frac1x,\qquad B=\sum_{i=1}^r\frac{b_i}{a_i^2}.\]
The level contains exactly $r^{k-1}$ distinct positive integers. Ordering them increasingly, the $j$th is at least $j$, so
\[0<S_k\le H_{r^{k-1}}\le1+(k-1)\log r.\tag{1}\]
Since all $a_i,b_i\ge1$, the minimum entry increases by at least one at each step. Thus every entry of $A_k$ is at least $k$.

The defining multiset recurrence, before any assumption of distinctness, gives
\[S_{k+1}=\sum_{x\in A_k}\sum_i\frac1{a_ix+b_i}.
\]
The exact identity
\[\frac1{a_ix+b_i}=\frac1{a_ix}-\frac{b_i}{a_ix(a_ix+b_i)}
\ge\frac1{a_ix}-\frac{b_i}{a_i^2x^2}
\]
therefore implies
\[S_{k+1}\ge C S_k-B\sum_{x\in A_k}\frac1{x^2}
\ge\left(C-\frac Bk\right)S_k.\tag{2}\]
Choose a fixed integer $K>B/(C-1)$ and set $\rho=C-B/K>1$. For every $k\ge K$, (2) gives $S_{k+1}\ge\rho S_k$, and consequently
\[S_k\ge\rho^{k-K}S_K.\tag{3}\]
The positive constant $S_K$ makes (3) exponential in $k$, contradicting (1). Some level must contain a repeated integer, proving the assertion.

\subsection*{The strict reciprocal threshold is meaningful}
For comparison, the two maps $x\mapsto2x+1$ and $x\mapsto2x+2$ have reciprocal sum exactly one. At depth $d$, their values at $1$ are
\[2^d+\sum_{j=0}^{d-1}2^j c_j,\qquad c_j\in\{1,2\}.
\]
After subtracting $\sum_{j=0}^{d-1}2^j$, the remaining digits $c_j-1$ are binary. Different words therefore yield different values at the same depth. Thus replacing $>1$ by $\ge1$ would make the statement false.

\subsection*{Outcome and attribution}
\textbf{SOLVED.} All estimates and the contradiction are elementary and included. The proof is the harmonic-weight method written by Kevin Barreto in the primary discussion, related there to Klarner's earlier affine-semigroup theorem; it is not a new solution. The stronger equality-of-affine-maps assertion is not required or claimed here.
Sources: \url{https://www.erdosproblems.com/481} and \url{https://www.erdosproblems.com/forum/discuss/481}, checked 24 September 2026. The existing Lean result belongs to its cited authors, and no local formal check of this text is claimed.

\subsection*{COMPLETION ESTIMATE}
\textbf{COMPLETION: 100\%}
